\section{Numerical Results}

\subsection{Simulation Setup}
We evaluate the performance of the proposed IB-CellFree-GNN in a simulated uplink Cell-Free Massive MIMO environment. The simulation area covers a $1 \times 1$ km$^2$ square with wrapped-around boundary conditions to approximate an infinite network. We deploy $L=16$ APs and $K=8$ single-antenna users (UEs), uniformly distributed within the area. Each AP is equipped with $N=4$ antennas unless otherwise specified. The channel model accounts for large-scale fading (path loss and shadowing) based on the 3GPP Urban Micro (UMi) model, and small-scale Rayleigh fading. The system bandwidth is 20 MHz, and the noise power density is $-174$ dBm/Hz. We compare the proposed method against three baselines: (1) \textbf{Centralized MMSE (C-MMSE)}, which assumes infinite fronthaul capacity and serves as a performance upper bound; (2) \textbf{Distributed Full-Precision (Dist-Full)}, where APs transmit 32-bit floating-point LMMSE estimates; and (3) \textbf{Distributed Fixed Quantization (Dist-Q2/Q4)}, where APs employ fixed 2-bit or 4-bit uniform quantization. The proposed QAT-GNN is trained to satisfy an average fronthaul constraint of $C_{target}=2.5$ bits per complex sample.

\subsection{BER Performance Analysis}
Fig. \ref{fig:ber_snr} illustrates the Bit Error Rate (BER) performance versus the transmit power $P_{tx}$ for the proposed scheme and the baselines. The results indicate that the proposed QAT-GNN significantly outperforms the fixed 2-bit quantization scheme (Dist-Q2) across the entire SNR regime. Specifically, at $P_{tx}=10$ dBm, QAT-GNN achieves a BER of approximately $1.1\%$, whereas Dist-Q2 saturates at $2.0\%$, representing a reduction in error rate of nearly $45\%$. Notably, the performance curve of QAT-GNN closely tracks the Distributed Full-Precision (Dist-Full) baseline, despite consuming an average of only 2.5 bits per sample compared to the 32 bits required for full precision. This efficiency stems from the Resource-Inference Module's ability to dynamically allocate higher bit-widths to APs with favorable channel conditions while suppressing noise-dominated signals from distant APs.

\subsection{Robustness to AP Dropouts}
To evaluate the resilience of the proposed architecture against fronthaul link failures, we simulate scenarios where a random subset of APs is disconnected from the CPU. Fig. \ref{fig:robustness} presents the BER performance under $0\%$, $25\%$, and $50\%$ AP dropout rates. The results demonstrate exceptional robustness: even when $50\%$ of the APs are disconnected (leaving only 8 active APs), the QAT-GNN maintains a BER of $2.1\%$ at $P_{tx}=20$ dBm. In contrast, conventional fixed quantization schemes typically degrade severely under such loss of spatial diversity. The Bit-Aware GNN at the CPU effectively compensates for the missing information by recalibrating the aggregation weights of the remaining active APs, leveraging the learned spatial correlation to mitigate the impact of the coverage holes.

\subsection{Scalability with User Load}
Fig. \ref{fig:scalability} analyzes the impact of user density on system performance by comparing a light load scenario ($K=4$) with a heavy load scenario ($K=12$). Under light loading, the gap between QAT-GNN and the fixed Dist-Q2 scheme is relatively small, as interference is manageable. However, as the user load increases to $K=12$, the performance gap widens significantly. At $P_{tx}=0$ dBm, QAT-GNN achieves a BER of $2.81\%$, while Dist-Q2 degrades to $8.60\%$. This substantial gain in the interference-limited regime confirms that the proposed task-oriented quantization does not merely reconstruct signals but actively preserves the discriminative features required for interference cancellation, which becomes critical as the number of users scales.

\subsection{Complexity and Antenna Scaling}
Finally, we investigate the trade-off between computational complexity and antenna array size in Fig. \ref{fig:antenna}. We compare the BER performance for APs equipped with $N=2$ and $N=8$ antennas. The results show that increasing $N$ from 2 to 8 yields a massive performance gain, reducing the BER by approximately two orders of magnitude at high SNR. Crucially, the computational overhead of the Policy Network remains constant regardless of $N$, as it operates on the compressed scalar statistics of the LMMSE output. A complexity analysis reveals that for $N=8$, the Policy Network adds only $\sim 6\%$ overhead to the local LMMSE processing. This confirms that the proposed method scales efficiently to massive MIMO configurations, leveraging large array gains without incurring prohibitive computational costs at the edge.